% Working outline: old prose cleared on 2026-11-09.
\section{Explored and set aside}
\label{sec:setaside}

\subsection{Four topic models that did not survive benchmarking}
\label{sec:setaside-models}

English and German use separate topic-discovery runs. \measured{} An earlier mixed-language run
produced mostly topics containing almost only English or almost only German documents.\footnote{
In the combined-language BERTopic runs, approximately 72--75\% of topics
had at least 95\% of their documents in one language. Noise assignments were excluded.
The diagnostic is saved in \texttt{output/stage2\_clean/}\allowbreak\texttt{lang\_purity\_report.json}.}
This motivated separate models.

\subsection{A separate cross-granularity matching diagnostic}
\label{sec:method-crossgran}

This exploratory diagnostic compares coarse LDA topics with fine BERTopic topics. It is
separate from the hierarchy used for RQ3 (Section~\ref{sec:rq3-hierarchy}) and leaves production
labels unchanged. For each fine topic $f$, its candidate parent $c$ is the coarse topic with
the largest document overlap. This adapts the pairing-strength idea from Bidirectional Topic
Matching~\cite[p.~4, Eq.~(1)]{adam2024btm}.

The script pairs the two models' assignments by page identifier. Overlap uses the resulting
rows carrying non-noise labels from both models; repeated rows are retained. It is the share
of those rows in $f$ whose coarse label is $c$.
The \emph{lead} is the largest overlap minus the second-largest overlap, with a missing
second share treated as zero. The candidate pair receives the weighted score

\begin{equation}
\mathrm{conf}(f, c) \;=\; 0.5 \cdot \mathrm{overlap}(f,c) \;+\; 0.3 \cdot \mathrm{sim}(f,c)
\;+\; 0.2 \cdot \mathrm{bidir}(f,c)
\end{equation}

where $\mathrm{sim}$ is the cosine similarity between the two topic centroids. For each coarse
topic, fine topics are ranked by overlap multiplied by their number of rows in the paired
assignments. The indicator $\mathrm{bidir}$ is one when $f$ ranks among the top three for $c$, and zero
otherwise. This is the contribution check used by the script.

As in TopicCheck~\cite[pp.~177--178]{chuang2015topiccheck}, a topic may remain without an
accepted match. This diagnostic applies its checks in the following order:
\begin{enumerate}[leftmargin=1.6em,itemsep=2pt]
\item \textsc{Reject} if overlap is below $0.35$ or similarity is below $0.25$.
\item Otherwise, mark \textsc{ambiguous} if the lead is below $0.25$.
\item Otherwise, mark \textsc{low confidence} if the weighted score is below $0.45$.
\item Otherwise, accept the candidate as a \textsc{match}.
\end{enumerate}

\begin{table}[h]\centering
\caption{Outcomes of the separate matching diagnostic at \texttt{min\_topic\_size=100}.
Each fine topic contributes one decision.}
\label{tab:crossgran}
\begin{tabular}{lrrrr}
\toprule
Language & \textsc{match} & \textsc{ambiguous} & \textsc{low conf.} & \textsc{reject} \\
\midrule
English & 99 & 12 & 0 & 1 \\
German  & 69 & 4  & 0 & 5 \\
\bottomrule
\end{tabular}
\end{table}

\measured{} The saved candidate-parent similarities range from $0.8927$ to $0.9993$ in English
and from $0.9300$ to $0.9991$ in German. These ranges describe the parent selected for each fine
topic. The saved files do not contain similarities for every possible fine--coarse pair.\footnote{\raggedright
Ranges, score minima and decision counts are recomputed from
\texttt{output/crossgran/}\allowbreak\texttt{crossgran\_\{eng,deu\}.json}.
Each file stores one record per fine topic.}

\measured{} No saved candidate failed the similarity check. Every rejection in
Table~\ref{tab:crossgran} was caused by insufficient overlap. Among candidates reaching the
weighted-score check, the lowest score was $0.5617$ in English and $0.4952$ in German.
Both exceed its cutoff, explaining the empty \textsc{low confidence} category. A
\textsc{match} label means only that the candidate passed these checks.

\subsection{Automatically generated keywords versus a human keyword list}
\label{sec:setaside-keywords}

\subsection{Search speed: a real engineering result with no research claim}
\label{sec:setaside-speed}

\subsection{A hyperlink topic graph}
\label{sec:setaside-graph}

\subsection{Category-filtered search, and why it cannot support a quality claim}
\label{sec:setaside-utility}

\subsection{Persona prompting the embedder}
\label{sec:setaside-persona}

A persona instruction asks the model to adopt a role. An earlier comparison used the role of
a research librarian. The motivation is given in Section~\ref{sec:method-embed}. The two
instructions were:

\begin{description}[leftmargin=1.3em,style=nextline,itemsep=3pt]
\item[Original.] ``Identify the main research topic of this web page''
\item[Librarian.] ``You are an expert research librarian. Identify the specific scientific
research theme of this web page.''
\end{description}

These were exploratory runs before the final production setup. Each language sample contained
50{,}000 entries, including some repeated pages. BERTopic used
\texttt{min\_topic\_size=100} throughout. Table~\ref{tab:prompt-comparison} reports the saved scores after
pages initially left outside any topic had been reassigned. Every sampled entry then received
a topic.\footnote{\raggedright Scores are from
\texttt{output/stage2\_lang/}\allowbreak\texttt{\{eng,deu\}/}\allowbreak\texttt{eval\_results.json}
and \texttt{output/stage2\_lang\_persona/}\allowbreak\texttt{\{eng,deu\}/}\allowbreak\texttt{eval\_results.json},
entry \texttt{BERTopic-min100-reduced}. Sample identities are stored in
\texttt{sample\_data/sample\_data.parquet} under each language directory.}
The scores are defined in Section~\ref{sec:method-metrics}. Word coherence used each run's
sample texts. Silhouette used 10{,}000 sampled entries after UMAP reduced the vectors to
50 coordinates.

\noindent\begin{minipage}{\linewidth}
\centering
\captionof{table}{Earlier instruction comparison after reassignment. \emph{Topics} gives the
number of discovered groups. Higher values are better for NPMI, $C_v$ and silhouette.}
\label{tab:prompt-comparison}
\begin{tabular}{llrrrr}
\toprule
Language & Instruction & Topics & NPMI & $C_v$ & Silhouette \\
\midrule
English & Original  & 110 & $-0.138$ & 0.429 & $-0.082$ \\
English & Librarian & 109 & $-0.143$ & 0.464 & $-0.037$ \\
German  & Original  &  77 & $-0.238$ & 0.402 & $-0.130$ \\
German  & Librarian &  70 & $-0.251$ & 0.352 & $-0.227$ \\
\bottomrule
\end{tabular}
\end{minipage}

\measured{} With the librarian instruction, English $C_v$ and silhouette are higher, while
NPMI is lower. In German, all three scores are lower. The changes therefore show no consistent
gain across the scores and languages. German coherence covers only 65 of the original run's
77 topics and 59 of the librarian run's 70 topics. The other topics had too few words in the
evaluation vocabulary, the list retained for coherence scoring.

The comparison has important limits. \measured{} The saved samples contain different pages
and include repeated entries. Comparing their page identifiers gives 7{,}809 shared distinct
English pages and 20{,}789 shared distinct German pages. The longer instruction also leaves
less room for page text in each chunk.
These differences prevent isolating the effect of instruction wording. This comparison covers
one alternative instruction and does not establish that other prompts would fail.

